\section{What was wrong}
\label{sec:wrong}

The problems are presented from the smallest to the most serious,
which is also the order in which the mechanization found them. Each is a
theorem checked by Rocq, and the theorem's name is given so that the
statement can be inspected. All of them are stated against a transcription of the original definitions.

\subsection{Errors that can be fixed in place}
\label{sec:local-errors}

Table~\ref{tab:local} lists six original claims that are false as stated.
None is deep: each is fixed by correcting the one definition or side
condition involved, and the rest of the design is unaffected.



\begin{table}[!ht]
\small
\begin{tabularx}{\textwidth}{@{}p{0.27\textwidth}X p{0.23\textwidth}@{}}
\toprule
Original claim & Counterexample & Repair \\
\midrule
\rn{assign} is sound for every invariant
  & With $J$ the assertion $x = 0$, the axiom derives
    $\spec{J}{\bexpt}{\cassign{x}{1}_1}{\bexpt}$, and the assignment
    falsifies $J$ with the lock free.
    \rocq{printed_assignment_invalid}
  & Require $x \notin \fv{J}$. \\[2pt]
Buffer overriding is associative (Prop.~2)
  & Three one-element buffers writing $1$, $2$, $3$ to one address: the
    list $[2,1,3]$ is in $a \override (b \override c)$ but not in
    $(a \override b) \override c$.
    \rocq{printed_buffer_override_not_associative}
  & A write from the later buffer may be taken only after the earlier
    buffer's writes to that address. \\[2pt]
Spatiotemporal separation is associative
  & Compatibility ignores blocked processors, so whether two memories
    compose can depend on a third that supplies the blocking fact.
    \rocq{printed_spatiotemporal_not_associative}
  & Check conflicts whether or not the processor is blocked. \\[2pt]
Assertions are closed under flushing (Prop.~9)
  & $\bexpf \fseqri \fbar{1}$ holds vacuously while processor $1$ has a
    pending write, and fails once the write flushes.
    \rocq{original_flushing_closure_false}
  & Implications quantify over all futures of the current memory. \\[2pt]
Conjecture~1
  & A cell duplicated across $\shash$ with equal values is not a spatial
    split. \rocq{printed_conjecture1_false}
  & Not needed by the final proof. \\[2pt]
\rn{unlock} is sound for every invariant
  & With $J = \bexpf$ and the lock held, $\cunlock_i$ is enabled and
    nothing restores $J$. \rocq{original_arbitrary_invariant_unlock_invalid}
  & See Section~\ref{sec:changes}. \\
\bottomrule
\end{tabularx}
\caption{Original claims that fail, each with a checked counterexample.}
\label{tab:local}
\end{table}

Two remarks. The repaired compatibility relation is stronger than the
original one, so it rejects some compositions the dissertation admits. Those
are states in which a blocked processor has a pending write to a cell that
someone else owns. And the repaired implications are
the ordinary Kripke reading of implication over the flushing preorder; the
three adjunctions between conjunctions and implications still hold.

The failure of \rn{unlock} points at something larger. The original safety
predicate asks for the invariant only while the lock is free. Once a
thread holds the lock, the predicate has forgotten that $J$ held at the
moment of acquisition, so neither \rn{unlock} nor \rn{inv} can be justified
from it; \rn{inv} on its own takes a true specification to a false one
(\rocq{original_inv_alone_counterexample}). That alone would call for a different predicate. The next
subsection shows that it calls for different rules.

\subsection{The original rules are inconsistent}
\label{sec:inconsistent}

The derivations below use the dissertation's own points-to, $x \fpointsto
v$, and are checked against its own interpretation of assertions.

\paragraph{History can be invented.}
A model is a memory and records nothing of how that memory came about. So a
current cell is also a model of any stale value followed by that cell:
\begin{equation}
  \label{eq:invent}
  x \fpointsto 7 \fsep \flock{i}
  \;\sentails\;
  x \fpointsto 5 \fseq (x \fpointsto 7 \fsep \flock{i}).
\end{equation}
The stale cell is overridden by the later one, and nothing remains of it.
This is the non-cancellativity that the dissertation notes in passing.

\paragraph{A $\cskip$ that changes memory.}
Let $J$ be $\femp \disj x \fpointsto 7$. The following is a derivation in
the original calculus.
\begin{enumerate}[nosep]
  \item $\spec{\femp}{J \fsep x \fpointsto 7 \fsep \flock{i}}{\cskip_i}
        {J \fsep \femp \fsep \flock{i}}$,
        by \rn{skip} and \rn{cons}: $J$ cannot be its second disjunct, since
        the cell would occur twice, and in the first disjunct the cell can
        be reassigned to $J$.
  \item $\spec{J}{x \fpointsto 7 \fsep \flock{i}}{\cskip_i}
        {\femp \fsep \flock{i}}$, by \rn{inv}. The thread has given its
        cell to the invariant.
  \item $\spec{J}{x \fpointsto 5 \fseq (x \fpointsto 7 \fsep \flock{i})}
        {\cskip_i}{x \fpointsto 5 \fseq (\femp \fsep \flock{i})}$,
        by \rn{frame-tm}.
  \item $\spec{J}{x \fpointsto 7 \fsep \flock{i}}{\cskip_i}
        {x \fpointsto 5 \fsep \flock{i}}$, by \rn{cons}, using
        \eqref{eq:invent} on the left and the unit law for $\femp$ on the
        right.
\end{enumerate}
The conclusion says that a thread holding the lock and a cell containing
$7$ may do nothing and then own that cell containing $5$. In the state
where processor $1$ holds the lock and $x$ contains $7$, the precondition
holds and the postcondition does not.
(\rocq{printed_points_stale_value_unsound}; the mechanized statement parses
\rn{inv}'s premise as $J \fsep (P \fsep \flock{i})$.)

\paragraph{A $\cskip$ that drains another processor's buffer.}
Spatial conjunction takes the union of the completeness sets and does not
require them to be disjoint, so an empty neighbour can copy a cell's marks:
$x \fpointsto 7 \fsep \femp \sentails x \fpointsto 7 \fsep \fbar{2}$. With
$J$ the cell $x \fpointsto 7$, \rn{skip}, \rn{cons} and \rn{inv} give
$\spec{J}{\femp \fsep \flock{1}}{\cskip_1}{\fbar{2} \fsep \flock{1}}$, and
\rn{frame-tm} with a pending write $R = y \fwrite{2} 17$ gives
\[
  \spec{J}{R \fseq (\femp \fsep \flock{1})}{\cskip_1}
          {R \fseq (\fbar{2} \fsep \flock{1})}.
\]
When processor $1$ holds the lock and processor $2$'s write to $y$ is still
buffered, the precondition holds. The postcondition places a barrier for
processor $2$ after its pending write, which that memory does not satisfy.
(\rocq{printed_points_foreign_barrier_unsound}.)

\paragraph{The strong temporal frame rule does not escape.}
One might hope that \rn{frame-stm} is safe, since its frame is also
spatially separate. It is not, because $R \fsseq P$ abbreviates
$(R \fsep P) \conj (R \fseq P)$ and the two conjuncts may be witnessed by
different splits. Take $R = \femp \disj x \fpointsto 5$. In the state above,
the spatial conjunct of $R \fsseq (x \fpointsto 7 \fsep \flock{1})$ holds
with $R$ empty, and the temporal conjunct holds with $R$ the stale cell.
Applying \rn{frame-stm} to step~2 yields a postcondition that the unchanged
state cannot satisfy. (\rocq{printed_strong_frame_counterexample}; this
derivation uses \rn{inv} and \rn{frame-stm} but not \rn{frame-tm}.)

All three conclusions are also invalid for the original safety predicate,
but that is the lesser point. They are false under the
dissertation's informal reading of a specification, and under any semantics
for which a true specification describes what the program does.

\subsection{Why no semantics could repair it}
\label{sec:trilemma}

The first derivation is an instance of a general conflict.

\begin{proposition}[\rocq{trilemma}]
\label{prop:trilemma}
Suppose a logic of triples has an unrestricted temporal frame rule, the
rule of consequence, and the entailment $R \fseq \femp \sentails R$.
Suppose $\htriple{P}{c}{\femp}$ is provable. Then $\htriple{R \fseq P}{c}{R}$
is provable for every $R$. Hence if some execution of $c$ starts in a state
satisfying $R \fseq P$ and ends in a state that does not satisfy $R$, the
logic is unsound.
\end{proposition}

The proof is two rule applications. The content is in the last sentence. A
triple $\htriple{P}{c}{\femp}$ with $P$ owning a cell is a \emph{transfer-out
contract}: after $c$, the cell belongs to someone else. Whenever temporal
conjunction lets later values override earlier ones, $R$ may describe an
older state of the same cell, as the store axiom's own postcondition does.
Framing $R$ across the transfer then hands the thread a stale copy of a
cell it no longer owns.

So four things cannot all be kept: the unrestricted temporal frame rule,
the right-unit law, transfer-out contracts ending in $\femp$, and
overlapping temporal history. A sound logic must give one up, and the
choice determines its character:
\begin{itemize}[nosep]
  \item \emph{No transfer out.} Cells may be taken from the invariant and
        never returned. This is sound and useless for a shared container.
  \item \emph{No right unit.} Let $\femp$ record that a cell was given
        away. Framing survives, and every postcondition that mentions
        $\femp$ changes meaning.
  \item \emph{Transfer leaves a mark.} Return an explicit tombstone
        assertion instead of $\femp$. The algebra survives; every later
        specification carries the tombstones.
  \item \emph{Restrict the frame.} Forbid temporal frames that overlap the
        cells being given away. Everything else is kept.
\end{itemize}
The first two were built and abandoned during this work, for the reasons
given. The repaired calculus takes the fourth.
Proposition~\ref{prop:trilemma} shows that some restriction of this kind is
needed, and the derivations above supply histories that do become false
after a transfer. It does not show that the particular restriction adopted
is the weakest possible.

The second derivation has a different cause and a different repair. The
all-processor completeness of points-to, and more generally the
completeness marks of whatever the invariant holds, should not pass to a
thread that acquires the lock. A thread can vouch only for its own buffer.

\subsection{The dissertation's own example}
\label{sec:own-example}

Section~4.1 of the dissertation proves a small message-passing program
correct, using an invariant that mentions two registers and the derived
rule \rn{atomic}. That proof does not carry over to the repaired calculus:
\rn{atomic} and \rn{share} in their original forms, and
auxiliary-variable elimination, are not available, and the concurrency rule now forbids either thread from
modifying a variable free in the invariant. The proof does not carry
over, although its conclusion is still true.
Section~\ref{sec:own-example-repaired} derives the same program from a
different invariant and recovers the original conclusion.
